% ---- pg-458 （书页 446） ----

\[
\mathbf{R}_{11}=\begin{bmatrix}1.000&0.785&0.810&0.775\\0.785&1.000&0.816&0.813\\0.810&0.816&1.000&0.845\\0.775&0.813&0.845&1.000\end{bmatrix}
\]
\[
{\small
\mathbf{R}_{12}=\mathbf{R}'_{21}=\begin{bmatrix}0.086&0.144&0.140&0.222&0.101&0.189&0.199&0.239\\0.200&0.119&0.211&0.301&0.223&0.221&0.274&0.235\\0.041&0.060&0.126&0.120&0.039&0.108&0.139&0.100\\0.228&0.122&0.277&0.214&0.201&0.156&0.271&0.171\end{bmatrix}
}
\]
\[
{\small
\mathbf{R}_{22}=\begin{bmatrix}1.000&0.562&0.457&0.579&0.802&0.595&0.512&0.492\\0.562&1.000&0.360&0.705&0.578&0.796&0.413&0.739\\0.457&0.360&1.000&0.273&0.606&0.337&0.798&0.240\\0.579&0.705&0.273&1.000&0.594&0.725&0.364&0.711\\0.802&0.578&0.606&0.594&1.000&0.605&0.698&0.605\\0.595&0.796&0.337&0.725&0.605&1.000&0.428&0.697\\0.512&0.413&0.798&0.364&0.698&0.428&1.000&0.394\\0.492&0.739&0.240&0.711&0.605&0.697&0.394&1.000\end{bmatrix}
}
\]

求样本典型变量及它们的相关系数. 解释这些量. 第一典型变量分别很好地综合了其各自变量组吗? 请解释.

\noindent 10.18\quad \hyperref[tab:7.7]{表 7.7} 中的数据含有纸浆纤维和纸的特性度量. 为与本章的符号相对应, 令纸的特性为
\[
x_1^{(1)}=\text{折断强度},\quad x_2^{(2)}=\text{弹性模量},\quad x_3^{(1)}=\text{失效应力},\quad x_4^{(1)}=\text{爆裂强度}
\]
纸浆纤维特性为
\[
x_1^{(2)}=\text{算术纤维长度},\quad x_2^{(2)}=\text{长纤维比例},\quad x_3^{(2)}=\text{细纤维比例},\quad x_4^{(2)}=\text{零伸展抗拉应力}
\]
确定样本典型变量与相互间的相关系数. 问: 第一典型变量是关联的变量的好的度量吗? 作出解释. 在显著性水平 $\alpha=0.05$ 下检验典型关系. 解释具有显著性的典型变量.

\noindent 10.19\quad 参见例 9.6 中的奥林匹克十项全能成绩的相关矩阵. 求速度项目(100 米跑、400 米跑、跳远)与臂力项目(铁饼、标枪、铅球)的典型相关关系. 提示: 标准化的跑步项目的符号改变了, 从而对所有项目得分越高越好.

\par\vspace{10pt}{\centering\heiti\zihao{4}参考文献\par}\vspace{4pt}

\begin{enumerate}[leftmargin=2.2em,itemsep=2pt]
\item Andrews, D.F., and A. M. Herzberg. \textit{Data}. New York: Springer-Verlag, 1985.
\item Bartlett, M.S. “Further Aspects of the Theory of Multiple Regression.” \textit{Proceedings of the Cambridge Philosophical Society}, \textbf{34} (1938), 33--40.
\item Bartlett, M.S. “A Note on Tests of Significance in Multivariate Analysis.” \textit{Proceedings of the Cambridge Philosophical Society}, \textbf{35} (1939), 180--185.
\item Dunham, R.B. “Reaction to Job Characteristics: Moderating Effects of the Organization.” \textit{Academy of Management Journal}, \textbf{20}, no. 1 (1977), 42--65.
\item Hotelling, H. “The Most Predictable Criterion.” \textit{Journal of Educational Psychology}, \textbf{26} (1935), 139--142.
\item Hotelling, H. “Relations between Two Sets of Variables.” \textit{Biometrika}, \textbf{28} (1936), 321--377.
\end{enumerate}
